% ---------------------------------------------------------------------------
\section{Material e Métodos}
% ---------------------------------------------------------------------------

Para operacionalizar a geração assistida, a metodologia adota um fluxo de trabalho em \textit{pipeline}, dividido em três etapas sequenciais: (1) Processamento da Linguagem, (2) Validação Determinística e (3) Construção e Testes, todas integradas através de um \textit{addon} em desenvolvimento para o Godot Engine.

\textbf{1. Processamento da Linguagem:} Inicialmente, um LLM recebe a requisição estrutural do desenvolvedor/jogador em linguagem natural (e.g., `crie uma masmorra com uma sala inicial, uma bifurcação onde um lado contém um tesouro e o outro lado demanda uma chave localizada no tesouro para acessar a sala do chefe''). O modelo é instruído via \textit{Prompt Engineering} a omitir cálculos espaciais (X, Y, Z) e exportar as conexões semânticas na forma de um grafo serializado em JSON. Essa matriz lógica reflete relações hierárquicas, ramificações e pré-requisitos mecânicos (portas trancadas vs. chaves).

\textbf{2. Validação Determinística:} O Godot intercepta o arquivo JSON e processa o grafo gerado através de um sistema baseado em regras clássicas. O algoritmo averígua condicionais estritas necessárias para uma fase funcional. As regras avaliam premissas topológicas, como: a existência imperativa de uma única sala inicial; a alcançabilidade da sala do chefe; a não alocação de chaves atrás da porta que elas abrem, tornando a sala inacessível; e a garantia de que as salas bônus (como tesouros) não estrangulem a rota principal de progresso. Caso o grafo apresente anomalias lógicas impeditivas, o sistema está sendo modelado para recusar o mapa e retroalimentar o LLM com as diretrizes do erro lógico.

\textbf{3. Construção do Mapa e Testes Espaciais:} Validado o grafo, o gerador espacial entra em vigor. Ele realiza a alocação procedural \textit{tile-based} ou via \textit{prefabs} de salas de combate, corredores e transições \cite{xu2025constraint}. O algoritmo determinístico calcula dinamicamente as posições espaciais com base no tamanho individual de cada malha tridimensional, impedindo sobreposições visuais e inconsistência física (\textit{clipping}). Subsequentemente, a \textit{engine} aplica ferramentas nativas, como nós de \textit{NavigationRegion3D} e projeção de \textit{RayCasts}, para atestar que o caminho principal possui uma malha de navegação (NavMesh) intacta, conectando fisicamente a entrada à saída.